\begin{pastexam}[label={px:nsysu-112-osds-02}]{中山資工所甲組 112 學年度 第 2 題}
（計算）

The pseudo code for the insertion sort is given as follows.
\begin{minted}[fontsize=\small]{text}
Read input elements into a[ ]; //stored in a[1], a[2], ..., a[n]
for i = 2 to n do
{
   j = i;
   while (j >= 2) and  (a[j] < a[j-1]) do
   {
      Swap(a[j], a[j-1]); // exchange a[j] and a[j-1]
      j = j-1;
   }
}
\end{minted}
\begin{enumerate}[label=(\alph*)]
\item (5\%) Suppose that the input elements are 3, 9, 6, 5, 8, 2. How many times (counts) are needed for
the execution of \texttt{Swap( )}?
\item (5\%) Suppose that the input is a permutation of $\{1,2,3,\ldots,n\}$. What permutation will make
\texttt{Swap( )} be executed the least times? How many times?
\item (5\%) Suppose that the input is a permutation of $\{1,2,3,\ldots,n\}$. What permutation will make
\texttt{Swap( )} be executed the most times? How many times?
\end{enumerate}
\end{pastexam}

\begin{pastsoln}
答案：(a) $9$ 次；(b) $1, 2, \ldots, n$，$0$ 次；(c) $n, n-1, \ldots, 1$，$n(n-1)/2$ 次。

題目的版本用相鄰交換：每一輪把 $a[j]$ 和左邊比它大的鄰居交換，一路換到左邊不比它大為止。 一次交換讓
兩個相鄰而且逆序的數變成順序對，其他數對的前後關係不變，所以剛好消掉一個逆序對。 這和資結書
第~1.3~節「挪動次數等於逆序對個數」是同一個論證：\texttt{Swap} 執行的次數等於輸入的逆序對個數。

(a) $9$ 次。 每一輪之後的陣列：
\begin{center}
\begin{tabular}{cccl}
\toprule
$i$ & 這一輪的 $a[i]$ & 交換次數 & 這一輪之後 \\
\midrule
2 & 9 & 0 & 3 9 6 5 8 2 \\
3 & 6 & 1 & 3 6 9 5 8 2 \\
4 & 5 & 2 & 3 5 6 9 8 2 \\
5 & 8 & 1 & 3 5 6 8 9 2 \\
6 & 2 & 5 & 2 3 5 6 8 9 \\
\midrule
合計 & & 9 & \\
\bottomrule
\end{tabular}
\end{center}
用逆序對驗算：$(3,2)$、$(9,6)$、$(9,5)$、$(9,8)$、$(9,2)$、$(6,5)$、$(6,2)$、$(5,2)$、$(8,2)$，共 $9$ 個。

(b) $1, 2, \ldots, n$（已經排好），$0$ 次。 交換次數等於逆序對個數，不會是負的；等於 $0$ 表示沒有任何
逆序對，也就是已經由小到大排好，而 $\{1, \ldots, n\}$ 排好的排列只有這一個。

(c) $n, n-1, \ldots, 1$（完全相反），$n(n-1)/2$ 次。 $n$ 個位置共有 $\binom{n}{2}$ 對，每一對最多貢獻
一個逆序對，所以交換次數最多 $n(n-1)/2$；要每一對都是逆序，排列只能是完全相反的那一個。 直接數也一樣：
第 $i$ 輪的 $a[i]$ 比左邊每個數都小，要換 $i - 1$ 次，$\sum_{i=2}^{n} (i - 1) = n(n-1)/2$。
\end{pastsoln}
